% TC22 — แยกคิวเจรจานัด / เรียงนัดใกล้ถึง
\begin{testscenario}{22}{4}
  \tcmeta{TC22}{CS 2568/01}{ระบบเว็บแอปพลิเคชัน ``NUCARE''}{นายสมชาย ใจดี}{วันที่ 4 พฤศจิกายน พ.ศ.\ 2569}{1.0}
  \tcname{แยกคิวเจรจานัดของที่ปรึกษา และการเรียงนัดใกล้ถึง}
  \tcpre{มีหน่วยทดสอบ \tccode{consultation-negotiate.test.ts} · รันด้วย \tccode{npx tsx --test src/components/advisor/consultation-negotiate.test.ts}}
  \tcdesc{ตรวจสอบการจัดเลนเจรจานัด (รอฉันตอบ / ยังไม่เสนอ / รอนิสิต) การเรียงลำดับ และการจับนัดภายใน 7 วัน สำหรับ inbox งานวันนี้}

  \tccase{1}{แยกเลนตาม waitingOn ของข้อเสนอ}{%
    \begin{tcsteps}
      \item เรียก \tccode{getNegotiateLane} กับข้อเสนอที่รอ ADVISOR / STUDENT / ไม่มีข้อเสนอ
    \end{tcsteps}}{%
    ได้ waiting\_advisor / waiting\_student / no\_offer ตามลำดับ}{%
    ผ่านตามหน่วยทดสอบ}{Pass}

  \tccase{2}{เรียง ACCEPTED ให้รออาจารย์ขึ้นก่อน}{%
    \begin{tcsteps}
      \item เรียก \tccode{sortAcceptedNegotiations} กับรายการปน 3 เลน
    \end{tcsteps}}{%
    ลำดับเป็นรอฉันตอบ → ยังไม่เสนอ → รอนิสิต}{%
    ผ่าน}{Pass}

  \tccase{3}{เรียง SCHEDULED ตามวันเวลานัด}{%
    \begin{tcsteps}
      \item เรียก \tccode{sortScheduledByAppointment}
    \end{tcsteps}}{%
    นัดที่ใกล้ถึงมาก่อน · ไม่มีวันเวลาอยู่ท้าย}{%
    ผ่าน}{Pass}

  \tccase{4}{นับเลนและนัดใกล้ถึงใน 7 วัน}{%
    \begin{tcsteps}
      \item เรียก \tccode{countNegotiateLanes} และ \tccode{isUpcomingAppointment}
    \end{tcsteps}}{%
    นับถูกต้อง · นัดใน horizon ถูกจับ และนัดไกลเกิน 7 วันไม่ถูกจับ}{%
    ผ่าน · ชุดทดสอบ 7 เคสรวม referral-active ผ่านทั้งหมด}{Pass}
\end{testscenario}
